\chapter{Análisis de Pequeña Señal y Parámetros del Amplificador}
\label{chap:pequena_senal}

En este capítulo se formula el circuito lineal equivalente de pequeña señal para la configuración en surtidor común,
deduciendo analíticamente las expresiones de ganancia de tensión, ganancia de corriente, impedancia
de entrada e impedancia de salida. Asimismo, se describen los fundamentos circuitales de los métodos de medición
empleados posteriormente en el laboratorio.

\section{Modelo Lineal de Pequeña Señal del MOSFET}

Para señales de alterna de reducida amplitud superpuestas al punto de reposo $Q$, el comportamiento no lineal del
MOSFET puede aproximarse con alta precisión mediante un cuadripolo linealizado gobernado por su transconductancia $g_m$
y su resistencia de salida drenador-surtidor $r_{ds}$.

En la figura~\ref{crkt:pequena_senal} se presenta el esquema circuital equivalente de pequeña señal del amplificador
completo.

\begin{figure}[!ht]
  \centering
  \begin{circuitikz}[scale=0.9, transform shape]
  % Terminal de entrada vi
  \draw (0,0) node[ground] {}
    to[open, v^>=$v_i$] (0,2.5)
    to[short, o-*] (1.5,2.5) coordinate (in_node);

  % Resistencia RG = R1 // R2 a masa de CA
  \draw (in_node) to[R, l=$R_G$, -] (1.5,0) node[ground] {};

  % Compuerta aislada (circuito abierto)
  \draw (in_node) to[short] (3.0,2.5) coordinate (gate_pt);
  \draw (3.0,0) node[ground] {} coordinate (source_pt);
  \draw (gate_pt) to[open, v^>=$v_{gs}$] (source_pt);

  % Fuente controlada gm * vgs orientada hacia abajo (desde drenador hacia masa)
  \draw (5.5,2.5) coordinate (drain_pt)
    to[cI, l_=$g_m v_{gs}$] (5.5,0) node[ground] {};

  % Resistencia RD y RL en paralelo a masa
  \draw (drain_pt) to[short, *-*] (7.0,2.5) coordinate (rd_node)
    to[R, l=$R_D$, -] (7.0,0) node[ground] {};

  \draw (rd_node) to[short, -*] (8.5,2.5) coordinate (rl_node)
    to[R, l=$R_L$, -] (8.5,0) node[ground] {};

  % Terminal de salida sobre el drenador
  \draw (rl_node) to[short, -o] (9.8,2.5) node[above] {$v_L$} coordinate (out_node);
  \draw (9.8,0) node[ground] {} to[open, v_>=$v_L$] (out_node);
\end{circuitikz}

  \caption{Modelo circuital equivalente de pequeña señal para el amplificador en surtidor común.}
  \label{crkt:pequena_senal}
\end{figure}

\subsection{Cálculo de la Transconductancia (\texorpdfstring{$g_m$}{gm})}
La transconductancia dinámica cuantifica la sensibilidad con la que la tensión compuerta-surtidor controla la corriente
de drenador en el punto de reposo:
\begin{equation}
  g_m = \left.\frac{\partial i_D}{\partial v_{GS}}\right|_{Q} = K \cdot (V_{GSQ} - V_{th}) = \frac{2 \cdot I_{DQ}}{V_{GSQ} - V_{th}} = \sqrt{2 \cdot K \cdot I_{DQ}}
  \label{eq:gm_formula}
\end{equation}

Evaluando con los valores del diseño analítico normalizado:
\begin{equation}
  g_m = \frac{2 \cdot 0.9727\mA}{0.1900\V} \approx 10.601\mS = 10.601 \times 10^{-3}\ \Omega^{-1}
  \label{eq:gm_calculada}
\end{equation}

\subsection{Resistencia de Salida Dinámica (\texorpdfstring{$r_{ds}$}{rds})}
La resistencia interna del canal en saturación proviene de la modulación de longitud de canal (efecto Early en MOSFETs):
\begin{equation}
  r_{ds} = \left(\left.\frac{\partial i_D}{\partial v_{DS}}\right|_{Q}\right)^{-1} = \frac{1}{\lambda \cdot I_{DQ}}
\end{equation}
Para el transistor discreto 2N7000, el parámetro de modulación $\lambda$ es sumamente reducido, resultando
$r_{ds} > 100\kohm$ (en el modelo SPICE del fabricante $r_{ds} \approx 2.4 \times 10^{11}\ohm \approx \infty$). Dado que
$r_{ds} \gg R_D = 6.8\kohm$, su conductancia en paralelo es despreciable a los efectos prácticos ($r_{ds} \parallel R_D \approx R_D$).

\section{Deducción Analítica de los Parámetros del Amplificador}

A partir de la topología de la figura~\ref{crkt:pequena_senal}, se deducen analíticamente las magnitudes del amplificador.

\subsection{Impedancia de Entrada (\texorpdfstring{$Z_i$}{Zi})}
Debido a la aislación galvánica de compuerta, la impedancia intrínseca de entrada del transistor tiende a infinito
($r_{\mathrm{gate}} \approx \infty$). Por consiguiente, la impedancia que la etapa presenta al generador de excitación
coincide con el equivalente en paralelo de la red de polarización de compuerta:
\begin{equation}
  Z_i = R_1 \parallel R_2 = R_{eq} = 101.833\kohm
  \label{eq:Zi_analitica}
\end{equation}

\subsection{Impedancia de Salida (\texorpdfstring{$Z_o$}{Zo})}
Para evaluar la impedancia de salida se anula la fuente de señal de entrada ($v_i = 0$), lo cual impone $v_{gs} = 0$.
Al extinguirse la tensión compuerta-surtidor, el generador dependiente de transconductancia se anula ($g_m v_{gs} = 0$),
actuando como un circuito abierto. La impedancia de salida vista desde los terminales de carga resulta:
\begin{equation}
  Z_o = R_D \parallel r_{ds} \approx R_D = 6.800\kohm
  \label{eq:Zo_analitica}
\end{equation}

\subsection{Ganancia de Tensión (\texorpdfstring{$A_v$}{Av})}
La corriente generada por el transistor inyecta una caída de tensión sobre la resistencia de carga dinámica de alterna:
\begin{equation}
  v_L = -(g_m \cdot v_{gs}) \cdot (R_D \parallel R_L)
  \label{eq:vl_pequena_senal}
\end{equation}
donde el signo negativo modela el sentido opuesto de la corriente respecto a la referencia de salida.

Dado que el surtidor se encuentra dinámicamente cortocircuitado a masa a través de $C_2$, la tensión compuerta-surtidor
coincide con la tensión de entrada ($v_{gs} = v_i$). El cociente define la ganancia de tensión de la etapa:
\begin{equation}
  A_v = \frac{v_L}{v_i} = -g_m \cdot (R_D \parallel R_L)
  \label{eq:Av_analitica}
\end{equation}

Evaluando con los valores normalizados ($g_m = 10.601\mS$ y $R_D \parallel R_L = 4.048\kohm$):
\begin{equation}
  A_v = -(10.601 \times 10^{-3}\ \Omega^{-1}) \cdot (4047.6\ohm) \approx -42.908 \approx -42.91
  \label{eq:Av_normalizada_val}
\end{equation}

\subsection{Ganancia de Corriente (\texorpdfstring{$A_i$}{Ai})}
La corriente de alterna entregada a la resistencia de carga es $i_L = v_L / R_L$, mientras que la corriente demandada
por el amplificador al generador es $i_i = v_i / Z_i$. La ganancia de corriente resulta:
\begin{equation}
  A_i = \frac{i_L}{i_i} = \frac{v_L / R_L}{v_i / Z_i} = A_v \cdot \frac{Z_i}{R_L}
  \label{eq:Ai_analitica}
\end{equation}

Evaluando con los parámetros normalizados:
\begin{equation}
  A_i = -42.908 \cdot \frac{101.833\kohm}{10.0\kohm} \approx -436.95
  \label{eq:Ai_normalizada_val}
\end{equation}

En el Cuadro~\ref{tab:resumen_pequena_senal} se resumen los parámetros de pequeña señal calculados.

\begin{table}[!ht]
  \centering
  \resizebox{0.88\textwidth}{!}{
  \begin{tabular}{|c|c|c|c|}
    \hline
    \textbf{Parámetro de Alterna} & \textbf{Teórico (Diseño)} & \textbf{Con R Normalizadas} & \textbf{Error Relativo} \\
    \hline
     $g_m\ [\mS]$ & $10.569$ & $10.601$ & $+0.30\%$ \\
     $Z_i\ [\kohm]$ & $100.000$ & $101.833$ & $+1.83\%$ \\
     $Z_o\ [\kohm]$ & $6.750$ & $6.800$ & $+0.74\%$ \\
     $A_v$ & $-42.591$ & $-42.908$ & $+0.74\%$ \\
     $A_i$ & $-425.91$ & $-436.95$ & $+2.59\%$ \\
    \hline
  \end{tabular}
  }
  \caption{Comparación de parámetros dinámicos de pequeña señal teóricos y con componentes normalizados.}
  \label{tab:resumen_pequena_senal}
\end{table}

\section{Fundamentos de los Métodos de Medición en Banco}

Para caracterizar con precisión los parámetros del amplificador en el laboratorio sin alterar el régimen de polarización
continua, se implementaron dos métodos circuitales indirectos:

\subsection{Método experimental para Medición de \texorpdfstring{$Z_i$}{Zi}}
En la figura~\ref{crkt:medicion_zi} se esquematiza el método de intercalación de un resistor serie calibrado
$R_s = 100\kohm$ entre el generador senoidal y la entrada del amplificador.

\begin{figure}[!ht]
  \centering
  \begin{circuitikz}[scale=0.9, transform shape]
  % Generador y resistencia serie Rs
  \draw (0,0) node[ground] {}
    to[sV, l=$v_s$] (0,2.5)
    to[R, l=$R_s\ (100\kohm)$, *-] (2.5,2.5) coordinate (vi_node)
    to[short, -o] (2.5,2.5);

  \draw (0,0) to[short] (2.5,0);
  \draw (0,2.5) to[short, *-o] (0,3.2) node[above] {CH1 ($v_s$)};
  \draw (vi_node) to[short, *-o] (2.5,3.2) node[above] {CH2 ($v_i$)};

  % Bloque amplificador Surtidor Comun
  \draw (3.0, -0.5) rectangle (6.2, 3.0);
  \node at (4.6, 1.25) [align=center] {Amplificador\\Surtidor Común\\(2N7000)};

  \draw (vi_node) to[short] (3.0, 2.5);
  \draw (2.5, 0) to[short] (3.0, 0);

  % Carga y medicion de salida
  \draw (6.2, 2.5) to[short, -*] (7.6, 2.5) coordinate (vl_node)
    to[R, l=$R_L\ (10\kohm)$] (7.6, 0) node[ground] {};
  \draw (6.2, 0) to[short] (7.6, 0);
  \draw (vl_node) to[short, -o] (8.5, 2.5) node[right] {$v_L$ (Salida)};
\end{circuitikz}

  \caption{Disposición experimental para la medición de la impedancia de entrada $Z_i$ por método del resistor serie.}
  \label{crkt:medicion_zi}
\end{figure}

Midiendo con el osciloscopio la tensión antes del resistor ($v_s$) y después del resistor ($v_i$), la corriente de
entrada resulta:
\begin{equation}
  i_i = \frac{v_s - v_i}{R_s}
  \label{eq:ii_medicion}
\end{equation}
A partir de la cual se calcula la impedancia de entrada experimental:
\begin{equation}
  Z_i = \frac{v_i}{i_i} = R_s \cdot \frac{v_i}{v_s - v_i}
  \label{eq:Zi_medicion}
\end{equation}

\subsection{Método experimental para Medición de \texorpdfstring{$Z_o$}{Zo} }
Para evaluar experimentalmente la impedancia de salida de la etapa, se utiliza el montaje de inyección ilustrado en la
figura~\ref{crkt:medicion_zo}.

\begin{figure}[!ht]
  \centering
  \begin{circuitikz}[scale=0.9, transform shape]
  % Bloque amplificador Surtidor Comun
  \draw (1.2, -0.5) rectangle (4.2, 3.0);
  \node at (2.7, 1.25) [align=center] {Amplificador\\Surtidor Común\\(2N7000)};

  % Entrada pasivada (cortocircuitada a masa de CA)
  \draw (0.0, 2.0) to[short] (1.2, 2.0);
  \draw (0.0, 0) to[short] (1.2, 0);
  \draw (0.0, 0) to[short] (0, 2);
  \draw (0.6, 0) node[ground] {};

  % Red de inyeccion en salida (nodo de drenador)
  \draw (4.2, 2.0) to[short, -*] (5.4, 2.0) coordinate (vo_node)
    to[R, l_=$R_s$] (7.4, 2.0)
    to[sV, l_=$v_s$] (7.4, 0) node[ground] {};
  \draw (4.2, 0) to[short] (7.4, 0);

  \draw (vo_node) to[short, *-o] (5.4, 3.0) node[above] {CH2 ($v_o$)};
  \draw (7.4, 2.0) to[short, *-o] (7.4, 3.0) node[above] {CH1 ($v_s$)};
\end{circuitikz}

  \caption{Esquema de inyección de señal para la medición de la impedancia de salida $Z_o$.}
  \label{crkt:medicion_zo}
\end{figure}

Se cortocircuita a masa la entrada de señal en alterna ($v_i = 0$) para anular cualquier excitación propia del
transistor, conservando intacta su polarización en continua. En el terminal de drenador se desacopla la resistencia de
carga $R_L$ y se conecta el generador de funciones en serie con un resistor testigo calibrado de $R_s = 6.68\kohm$.

Midiendo con el osciloscopio la tensión del generador ($v_s$) en el canal CH1 y la tensión en el drenador ($v_o$) en el
canal CH2, la corriente inyectada sobre el nodo de salida es:
\begin{equation}
  i_o = \frac{v_s - v_o}{R_s}
  \label{eq:io_medicion}
\end{equation}
A partir de la cual la impedancia de salida queda determinada como:
\begin{equation}
  Z_o = \frac{v_o}{i_o} = R_s \cdot \frac{v_o}{v_s - v_o}
  \label{eq:Zo_medicion}
\end{equation}
